\documentclass[a4paper,12pt,fleqn]{article}

\setlength{\textwidth}{15.0cm}
\setlength{\textheight}{24.0cm}
\setlength{\topmargin}{0cm}
\setlength{\headsep}{0cm}
\setlength{\headheight}{0cm}
\pagestyle{plain}

\usepackage[dvips]{epsfig}
\usepackage{tikz}
\usepackage[margin=1in]{geometry}
\usepackage{amsmath,amssymb,amsthm}
\usepackage{comment}

\newtheoremstyle{break}{4pt}{4pt}{}{}{\bfseries}{\vspace{2 pt}}{\newline}{}
\theoremstyle{break}
\newtheorem{theorem}{Theorem}[subsection]
\newtheorem{example}[theorem]{Example}
\newtheorem{definition}[theorem]{Definition}
\newtheorem{notation}[theorem]{Notation}
\newtheorem{lemma}[theorem]{Lemma}
\newtheorem{conjecture}[theorem]{Conjecture}
\newtheorem{corollary}[theorem]{Corollary}
\newtheorem{technique}[theorem]{Technique}
\newtheorem{pythonn}[theorem]{Python Code}
\newtheorem{related}[theorem]{Related Work}

\begin{document}

\begin{comment}
\title{Unbiased Monte Carlo for Recursive Integrals}
\author{Isidoor Pinillo Esquivel}
\maketitle
\end{comment}

\begin{titlepage}
    \begin{center}
        \resizebox{3cm}{!}{\includegraphics{./vert2_kl_01.eps}}
        \ \
        \ \\
        \ \\
        \ \\
        \ \\
        \ \\
        \ \\
        \ \\
        \ \\
        \ \\
        \ \\
        \ \\
        \Large{M{\sc asterproef scriptie}}
        \ \\
        \ \\
        \ \\
        \huge{\bf{\em Unbiased Monte Carlo for Recursive Integrals}}
        \ \\
        \ \\
        \ \\
        \ \\
        \ \\
        \ \\
        \normalsize
        Auteur: {\em Isidoor Pinillo Esquivel}\\
        \ \\
        \ \\
        Promotor: {\em Wim Vanroose}\\
        \ \\
        \ \\
        \ \\
        \ \\
        \ \\
        \ \\
        \ \\
        A{\sc cademiejaar 2022-2023}

    \end{center}
\end{titlepage}




\newpage
\tableofcontents
\newpage

\begin{abstract}
    We will write this at the end.
\end{abstract}

\section{Introduction}

\subsection{Introductory Example}
Follows $y'=y$ example. This introduces recursive Monte Carlo and
highlight some advantages and problems of that example.

\subsection{Contributions}
We will write this at the end.

\subsection{Related Work}
work on
\begin{itemize}
    \item alternative methods for recursive integrals
    \item MC work on ODEs
    \item MC work on PDEs
    \item WoS
\end{itemize}
This is just to give a general overview we probably reference specific ideas when we first introduce them.

\section{Background}
\subsection{Modifying Monte Carlo}
Introduces Russian roulette, splitting, control variates, importance sampling and maybe quasi Monte Carlo with the
$y'=y$ example.

\subsection{Monte Carlo Trapezoidal Rule}
comparing normal vs Monte Carlo trapezoidal rule and highlighting the "half variance phenomenon".
+ maybe integrating polynomials for intuition

\subsection{Unbiased Non-Linearity}
\begin{example}[$y'=y^{2}$]
    see python note book
\end{example}

\begin{example}[$e^{E[X]}$]
    see python note book
\end{example}

\subsection{Recursion}

\begin{example}[coupled recursion]
    example with $y'=y$ (I need to redo this example)
\end{example}

\begin{example}[recursion in recursion]
    maybe induction in induction proof example
\end{example}

\begin{example}[tail recursion]
    discuss problems with implementing recursion and solutions.  \\
    inverse problem example
\end{example}

\subsection{Green Functions}
green function stuff that we will be needing, we aren't sure in how much detail we're going to go.

\begin{example}[numerical green functions]
    There will be probably some green functions that we need
    that don't have an analytic expression yet.
\end{example}

\section{$1$-Dimensional Recursive Integrals}

\subsection{Linear Recursive Integrals}
We have algo in mind for this case based on coupled recursion on disjunct sets.


\subsection{IVPs ODEs}
An IVP example probably using DRRMC maybe compare it to parareal. Maybe also non-linear algo

\subsection{BVPs ODEs}
A BVP example using yet another algo that hopefully has the half variance phenomenon.

\section{Higher Dimensional Recursive Integrals}
\subsection{Complicated Geometry}
\begin{example}[nasty $2$D integral]
    $2$D integral that is difficult because of its geometry
\end{example}

\subsection{Recursive Brownian Motion}
WoS like way to simulate Brownian motion which is related to the green function
of the heat equation

\begin{example}[recursive Brownian motion]
    see period5
\end{example}

\subsection{Heat Equation}
a geometric robust way to solve the heat equation and maybe a higher order method to solve
the heat equation

\subsection{Wave Equation}
probably won't get to it

\section{Appendix}
Derivation of the green functions and some expressions.

\end{document}
